
\section{二阶导数引理的带（多）参数变体}
\label{sec:variants-parameters}

在本节中，我们给出前面用过的、\trFloer{}算子二阶导数的诸引理的几种带（多）参数的版本。

为了证明\cref{thm:gamma-newton-picard}，我们需要\trFloer{}算子二阶导数的\cref{thm:df-drho-estimate}的一个变体：我们需要顾及 $H$ 与 $J$ 依赖于\trGluing{}所用的参数 $\rho$ 这一事实。这就是我们现在要做的事。

为了验证\cref{thm:composition-rho-derivative-estimate}的诸条件，我们证明了：存在一个常数 $C_{1}$，使得对一切满足 $\|X(s,t)\|\le r_{1}$ 的 $s$、$t$ 以及每个 $\rho$，
\[
\Bigl\|\frac{\partial\widetilde{\scr{F}}_{\rho}}{\partial\rho}(X)-\frac{\partial\widetilde{\scr{F}}_{\rho}}{\partial\rho}(0)\Bigr\|\le C_{1}\Bigl(\|X\|+\Bigl\|\frac{\partial X}{\partial s}\Bigr\|+\Bigl\|\frac{\partial X}{\partial t}\Bigr\|\Bigr)\text{。}
\]
计算导数 $\dfrac{\partial\widetilde{\scr{F}}_{\rho}}{\partial\rho}(X)$ 的公式，由于 $J$ 与 $H$ 对 $\rho$ 的依赖而含有两个新的项；它们是
\[
P(X)=\frac{\partial J}{\partial\rho}(s+\rho)\left(D_{1}\widetilde{\exp}_{(w_{\rho},X)}\left(\frac{\partial w_{\rho}}{\partial t}\right)+D_{2}\widetilde{\exp}_{(w_{\rho},X)}\left(\frac{\partial X}{\partial t}\right)\right)
\]
与
\[
Q(X)=\operatorname{grad}\frac{\partial H_{s+\rho,t}}{\partial\rho}\left(\widetilde{\exp}(w_{\rho},X)\right)\text{。}
\]
\cref{thm:a-b-estimates} 给出
\[
\|P(X)-P(0)\|\le C\Bigl(\|X\|+\Bigl\|\frac{\partial X}{\partial t}\Bigr\|\Bigr)\text{。}
\]
其次，对每个对子 $(s,t)$，
\[
\bigl(w_{\rho}(s,t),\sigma X(s,t)\bigr)\in W\times\overline{B}(0,r)\quad\text{（一个\trCompact{}集）}
\]
而
\[
\frac{\partial H_{s+\rho}}{\partial\rho}=0\qquad\text{对 } |s+\rho|\ge R\text{。}
\]
于是，
\[
\Bigl\|\Bigl(d\operatorname{grad}\frac{\partial H_{s+\rho,t}}{\partial\rho}\Bigr)_{\widetilde{\exp}(p,X)}\Bigr\|
\]
在\trCompact{}集 $[-R,R]\times S^{1}\times W\times\overline{B}(0,r)$ 上\trBounded，从而最终
\[
\begin{aligned}
\|Q(X)-Q(0)\|&\le\sup_{\sigma\in[0,1]}\Bigl\|\Bigl(d\operatorname{grad}\frac{\partial H_{s+\rho,t}}{\partial\rho}\Bigr)_{\widetilde{\exp}(w_{\rho},\sigma X)}\cdot D_{2}\widetilde{\exp}_{(w_{\rho},\sigma X)}\Bigr\|\\
&\le C\,\|X\|\text{，}
\end{aligned}
\]
这使得我们能够找到所要的估计并结束证明。

我们还将证明我们在\cref{subsec:construction-phi-lambda}中需要的那条\cref{thm:dn-op-estimate}的变体，它如下。

\emph{引理（\cref{thm:dfbar-lipschitz}）。}设 $r_{0}>0$。存在一个常数 $K>0$，使得对每个 $\rho\ge\rho_{0}$ 与每个 $(\lambda,Z)\in\RR\times W^{1,p}(\RR\times S^{1};\RR^{2n})$（其中 $\|Z\|\le r_{0}$），我们有
\[
\bigl\|(d\overline{F}_{\rho})_{(\lambda,Z)}-(d\overline{F}_{\rho})_{(0,0)}\bigr\|^{\mathrm{op}}\le K\bigl(\|Z\|_{W^{1,p}}+|\lambda|\bigr)\text{。}
\]

\begin{proof}
我们有
\[
\bigl\|(d\overline{F}_{\rho})_{(\lambda,Z)}-(d\overline{F}_{\rho})_{(0,0)}\bigr\|^{\mathrm{op}}\le\bigl\|(d\overline{F}_{\rho})_{(\lambda,Z)}-(d\overline{F}_{\rho})_{(\lambda,0)}\bigr\|^{\mathrm{op}}+\bigl\|(d\overline{F}_{\rho})_{(\lambda,0)}-(d\overline{F}_{\rho})_{(0,0)}\bigr\|^{\mathrm{op}}\text{。}
\]
让我们分别分析右端的两项。我们从第一项开始。设 $(a,Y)\in\RR\times W^{1,p}(\RR\times S^{1};\RR^{2n})$。我们有\footnote{译注：下一式取值点原书印作 $\sigma=\lambda$；按被导函数为 $\sigma\mapsto\scr{F}_{\rho}^{\Gamma_{\lambda+\lambda_{\star}+a\sigma}}\bigl(\exp_{w_{\rho}}(Z)\bigr)$，于 $\sigma=\lambda$ 处取值将落到参数 $\lambda+\lambda_{\star}+a\lambda$ 上，而下一行所得的 $a\Bigl(\frac{\partial J}{\partial\lambda}(\lambda+\lambda_{\star},\cdots)\cdots+\operatorname{grad}\frac{\partial H}{\partial\lambda}\Bigr)$ 恰是 $\sigma=0$ 处的链式结果，疑应于 $\sigma=0$ 处取值。此处径改；下一行诸式不受影响。——译者注}
\[
\begin{aligned}
(d\overline{F}_{\rho})_{(\lambda,Z)}(a,Y)&=(d\overline{F}_{\rho})_{(\lambda,Z)}(0,Y)+(d\overline{F}_{\rho})_{(\lambda,Z)}(a,0)\\
&=(dF_{\rho}^{\Gamma_{\lambda+\lambda_{\star}}})_{Z}(Y)+\frac{\partial}{\partial\sigma}\Bigm|_{\sigma=0}\Bigl(\scr{F}_{\rho}^{\Gamma_{\lambda+\lambda_{\star}+a\sigma}}\bigl(\exp_{w_{\rho}}(Z)\bigr)\Bigr)_{(Z_{i}^{\rho,Z})}\\
&=(dF_{\rho}^{\Gamma_{\lambda+\lambda_{\star}}})_{Z}(Y)+a\left(\frac{\partial J}{\partial\lambda}(\lambda+\lambda_{\star},s+\rho)\,\Bigl(\frac{\partial(\exp_{w_{\rho}}Z)}{\partial t}\right.\\
&\qquad\qquad\qquad\qquad\left.-X_{H_{s+\rho,t}^{\lambda+\lambda_{\star}}}\bigl(\exp_{w_{\rho}}(Z)\bigr)\Bigr)\right.\\
&\qquad\qquad\qquad\qquad\left.+\operatorname{grad}_{\exp_{w_{\rho}}(Z)}\frac{\partial H}{\partial\lambda}(\lambda+\lambda_{\star},s+\rho,t)\right)_{(Z_{i}^{\rho,Z})}
\end{aligned}
\]
（在这些表达式里，$(Z_{i}^{\rho,Z})$ 表示点 $\exp_{w_{\rho}(s,t)}Z(s,t)$ 处的那个\trOrthonormal\trFrame，它由\trFrame{} $(Z_{i}^{\rho}(s,t))$ 沿曲线 $r\mapsto\exp_{w_{\rho}(s,t)}(rZ(s,t))$ 作\trParallelTransport{}而得到）。\footnote{译注：本式、下文 $V_{\rho}(\lambda+\lambda_{\star},0)$ 的定义式与再下文 $\widetilde{V}_{\rho}(\lambda+\lambda_{\star},Z)$ 的定义式，原书均把参数方向场印作 $\frac{\partial J}{\partial\lambda}(\lambda+\lambda_{\star},s+\rho)\,\frac{\partial(\cdots)}{\partial t}+\operatorname{grad}_{(\cdot)}\frac{\partial H}{\partial\lambda}$ 之形。按本书约定（第~5~章）$\operatorname{grad}^{J}H=-JX_{H}$，而 $J$ 亦随 $\lambda$ 变动，对 $\lambda$ 求导尚应产生 $-\frac{\partial J}{\partial\lambda}X_{H}$ 一项，即参数方向场应为 $\frac{\partial J}{\partial\lambda}\Bigl(\frac{\partial(\cdots)}{\partial t}-X_{H}\Bigr)+\operatorname{grad}\frac{\partial H}{\partial\lambda}$（$X_{H}$ 的下标随该处的\trHamiltonian{}而定），原书疑漏 $-\frac{\partial J}{\partial\lambda}X_{H}$ 项。此疏漏与第~11~章的同一处（彼处已立总注径改）同源。三处径改。——译者注}

我们采用记号
\[
d\overline{F}_{\rho}(\lambda,Z)(a,0)=a\,\bigl(V_{\rho}(\lambda+\lambda_{\star},Z)\bigr)_{(Z_{i}^{\rho,Z})}
\]
并且注意到 $V_{\rho}(\lambda_{\star},0)$ 正是出现在计算 $L_{\rho}^{\Gamma}(a,0)$ 的公式中的那个\trVectorField{} $V_{\rho}$。

为了从上控制 $\bigl\|(d\overline{F}_{\rho})_{(\lambda,Z)}-(d\overline{F}_{\rho})_{(\lambda,0)}\bigr\|^{\mathrm{op}}$，我们写
\[
\begin{aligned}
&\bigl\|(d\overline{F}_{\rho})_{(\lambda,Z)}(a,Y)-(d\overline{F}_{\rho})_{(\lambda,0)}(a,Y)\bigr\|_{L^{p}}\\
&\qquad\le\bigl\|(dF_{\rho}^{\Gamma_{\lambda+\lambda_{\star}}})_{Z}(Y)-(dF_{\rho}^{\Gamma_{\lambda+\lambda_{\star}}})_{0}(Y)\bigr\|_{L^{p}}\\
&\qquad\qquad+|a|\,\Bigl\|\bigl(V_{\rho}(\lambda+\lambda_{\star},Z)\bigr)_{(Z_{i}^{\rho,Z})}-\bigl(V_{\rho}(\lambda+\lambda_{\star},0)\bigr)_{(Z_{i}^{\rho,Z})}\Bigr\|_{L^{p}}\text{。}
\end{aligned}
\]
为了控制这个和的第一项，我们使用\cref{thm:dn-op-estimate}，或者更精确地说，使用它对由\trHomotopy{} $\Gamma_{\lambda+\lambda_{\star}}$ 所定义的偶对 $(H,J)$ 的类似物——在\cref{thm:gamma-newton-picard}的证明中，我们已经提过这种类比的可能性。我们有
\[
\bigl\|(dF_{\rho}^{\Gamma_{\lambda+\lambda_{\star}}})_{Z}(Y)-(dF_{\rho}^{\Gamma_{\lambda+\lambda_{\star}}})_{0}(Y)\bigr\|_{L^{p}}\le k_{1}\,\|Z\|_{W^{1,p}}\|Y\|_{W^{1,p}}
\]
对一个不依赖于 $\rho$ 的常数 $k_{1}>0$ 成立。为了控制余下的那一项，我们把 $V_{\rho}(\lambda+\lambda_{\star},\bullet)$ 视为一个映射
\[
W^{1,p}(\RR\times S^{1};\RR^{m})\longrightarrow L^{p}(\RR\times S^{1};\RR^{m})\text{，}
\]
并（采用第~11~章的记号）写出
\[
\begin{aligned}
\widetilde{V}_{\rho}(\lambda+\lambda_{\star},Z)&=\frac{\partial\widetilde{J}}{\partial\lambda}(\lambda+\lambda_{\star},s+\rho)\left(\frac{\partial(\exp_{w_{\rho}}Z)}{\partial t}\right.\\
&\qquad\qquad\left.-X_{H_{s+\rho,t}^{\lambda+\lambda_{\star}}}\bigl(\exp_{w_{\rho}}(Z)\bigr)\right)\\
&\qquad+\operatorname{grad}_{\widetilde{\exp}_{w_{\rho}}(Z)}\left(\frac{\partial H}{\partial\lambda}(\lambda+\lambda_{\star},s+\rho,t)\right)\\
&=\frac{\partial\widetilde{J}}{\partial\lambda}(\lambda+\lambda_{\star},s+\rho)\left(D_{1}\widetilde{\exp}_{(w_{\rho},Z)}\left(\frac{\partial w_{\rho}}{\partial t}\right)+D_{2}\widetilde{\exp}_{(w_{\rho},Z)}\left(\frac{\partial Z}{\partial t}\right)\right.\\
&\qquad\qquad\qquad\left.-X_{H_{s+\rho,t}^{\lambda+\lambda_{\star}}}\bigl(\exp_{w_{\rho}}(Z)\bigr)\right)\\
&\qquad\qquad+\operatorname{grad}_{\widetilde{\exp}(w_{\rho},Z)}\,\frac{\partial H}{\partial\lambda}(\lambda+\lambda_{\star},s+\rho,t)\text{。}
\end{aligned}
\]

我们使用下面这个结果。

\begin{lemma}\label{thm:vtilderho-estimate}
设 $r_{0}>0$。存在一个常数 $K>0$，使得对每个 $\rho\ge\rho_{0}$ 与每个 $Z\in\overline{B}(0,r_{0})\subset W^{1,p}(\RR\times S^{1};\RR^{m})$，我们有
\[
\bigl\|\widetilde{V}_{\rho}(\lambda+\lambda_{\star},Z)-\widetilde{V}_{\rho}(\lambda+\lambda_{\star},0)\bigr\|_{L^{p}}\le K\,\|Z\|_{W^{1,p}}\text{。}
\]
\end{lemma}

\begin{proof}
让我们写
\[
\begin{aligned}
&\bigl\|\widetilde{V}_{\rho}(\lambda+\lambda_{\star},Z)-\widetilde{V}_{\rho}(\lambda+\lambda_{\star},0)\bigr\|_{L^{p}}\\
&\qquad\le\Bigl\|\frac{\partial\widetilde{J}}{\partial\lambda}(\lambda+\lambda_{\star},s+\rho)\Bigl[D_{1}\widetilde{\exp}_{(w_{\rho},Z)}\Bigl(\frac{\partial w_{\rho}}{\partial t}\Bigr)+D_{2}\widetilde{\exp}_{(w_{\rho},Z)}\Bigl(\frac{\partial Z}{\partial t}\Bigr)\\
&\qquad\qquad\qquad\qquad\qquad\qquad-D_{1}\widetilde{\exp}_{(w_{\rho},0)}\Bigl(\frac{\partial w_{\rho}}{\partial t}\Bigr)\Bigr]\Bigr\|_{L^{p}}\\
&\qquad+\Bigl\|\operatorname{grad}_{\widetilde{\exp}_{(w_{\rho},Z)}}\frac{\partial H}{\partial\lambda}(\lambda+\lambda_{\star},s+\rho,t)-\operatorname{grad}_{\widetilde{\exp}_{(w_{\rho},0)}}\frac{\partial H}{\partial\lambda}(\lambda+\lambda_{\star},s+\rho,t)\Bigr\|_{L^{p}}\text{。}
\end{aligned}
\]
映射
\[
\frac{\partial\widetilde{J}}{\partial\lambda}:[0,1]\times\RR\longrightarrow\operatorname{End}(TW)
\]
具有紧\trSupport。于是存在一个常数 $C>0$，使得
\[
\Bigl\|\frac{\partial\widetilde{J}}{\partial\lambda}(\lambda,s)\Bigr\|^{\mathrm{op}}\le C\qquad\text{对每个 }(\lambda,s)\in[0,1]\times\RR\text{。}
\]
于是，我们不等式右端的第一项可以从上控制为
\[
C\Bigl\|D_{1}\widetilde{\exp}_{(w_{\rho},Z)}\Bigl(\frac{\partial w_{\rho}}{\partial t}\Bigr)+D_{2}\widetilde{\exp}_{(w_{\rho},Z)}\Bigl(\frac{\partial Z}{\partial t}\Bigr)-D_{1}\widetilde{\exp}_{(w_{\rho},0)}\Bigl(\frac{\partial w_{\rho}}{\partial t}\Bigr)\Bigr\|_{L^{p}}\text{。}
\]
我们把\cref{thm:a-b-estimates}与\cref{thm:w1p-product-lp}应用于
\[
A(p,X,V)=D_{1}\widetilde{\exp}_{(p,X)}V\qquad\text{与}\qquad A(p,X,V)=D_{2}\widetilde{\exp}_{(p,X)}V\text{，}
\]
（采用与\cref{sec:proof-second}中相同的论证），以便把这个表达式从上控制为 $K\|Z\|_{W^{1,p}}$。

为了控制第二项（即带\trGradient{}的那一项），注意到 $Z$ 的 $W^{1,p}$ 上界 $\|Z\|_{W^{1,p}}\le r_{0}$ 蕴含 $Z$ 是\trBounded{}的：$\|Z\|_{L^{\infty}}\le r_{1}$（对某个 $r_{1}>0$）。考虑映射
\[
\begin{aligned}
\varphi:[0,1]\times\RR\times S^{1}\times W\times\overline{B}(0,r_{1})&\longrightarrow\RR^{m}\\
(\lambda,s,t,p,X)&\longmapsto\operatorname{grad}_{\widetilde{\exp}_{(p,X)}}\frac{\partial H}{\partial\lambda}(\lambda,s,t)\text{。}
\end{aligned}
\]
它是 $\cC^{1}$ 类的，并且具有紧\trSupport。于是存在一个常数 $C>0$，使得 $\|d\varphi\|^{\mathrm{op}}\le C$。由此推出，对\trEuclidean{}范数，
\[
\bigl\|\varphi(\lambda+\lambda_{\star},s+\rho,t,w_{\rho}(s,t),Z)-\varphi(\lambda+\lambda_{\star},s+\rho,t,w_{\rho}(s,t),0)\bigr\|\le C\,\|Z\|\text{。}
\]
取 $L^{p}$ 范数，我们便推出：我们不等式右端的第二项正如所愿地从上得到了控制：
\[
\begin{aligned}
\Bigl\|\operatorname{grad}_{\widetilde{\exp}_{(w_{\rho},Z)}}\frac{\partial H}{\partial\lambda}(\lambda+\lambda_{\star},s+\rho,t)-\operatorname{grad}_{\widetilde{\exp}_{(w_{\rho},0)}}\frac{\partial H}{\partial\lambda}(\lambda+\lambda_{\star},s+\rho,t)\Bigr\|_{L^{p}}\\
\le C\,\|Z\|_{L^{p}}\le C\,\|Z\|_{W^{1,p}}\text{。}
\end{aligned}
\]
这就结束了\cref{thm:vtilderho-estimate}的证明。
\end{proof}

但\cref{thm:dfbar-lipschitz}的证明还没有结束。正如我们在第~13~章中所做的那样，我们使用一幅交换图。
\[
\begin{tikzcd}[column sep=3.2cm,row sep=1.6cm]
W^{1,p}(\RR\times S^{1};\RR^{2n}) \ar[r,"{z\mapsto V_{\rho}(\lambda+\lambda_{\star},Z)}"] \ar[d,"i_{\rho}"']
& L^{p}(\RR\times S^{1};\RR^{2n}) \\
W^{1,p}(\RR\times S^{1};\RR^{m}) \ar[r,"{(\operatorname{Id},\widetilde{V}_{\rho})}"']
& W^{1,p}(\RR\times S^{1};\RR^{m})\times L^{p}(\RR\times S^{1};\RR^{m}) \ar[u,"p_{\rho}"]
\end{tikzcd}
\]

设 $\varphi=i_{\rho}$、$\psi=(\operatorname{Id},\widetilde{V}_{\rho})$ 与 $\chi=p_{\rho}$。为了完成所要上界的建立，我们需要下面的引理。

\begin{lemma}\label{thm:chi-psi-estimate}
存在一个常数 $k_{2}>0$（它不依赖于 $\rho$），使得当 $\|Z\|_{W^{1,p}}\le r_{0}$ 时，我们有
\[
\|\chi\psi\varphi(Z)-\chi\psi\varphi(0)\|_{L^{p}}\le k_{2}\,\|Z\|_{W^{1,p}}\text{。}
\]
\end{lemma}

\begin{proof}
映射 $\varphi$ 是线性的，它的范数关于 $\rho$ 一致\trBounded，于是
\[
\|\varphi(Z)\|_{W^{1,p}}\le Cr_{0}\text{。}
\]
由\cref{thm:vtilderho-estimate}，存在一个 $K>0$，使得
\[
\begin{aligned}
\|\psi\varphi(Z)-\psi\varphi(0)\|_{W^{1,p}\times L^{p}}&=\|\varphi(Z)-\varphi(0)\|_{W^{1,p}}+\bigl\|\widetilde{V}_{\rho}(Z)-\widetilde{V}_{\rho}(0)\bigr\|_{L^{p}}\\
&\le C\,\|Z\|_{W^{1,p}}+K\,\|Z\|_{W^{1,p}}\text{。}
\end{aligned}
\]
其次，我们令 $a=\psi(\varphi(Z))$ 与 $b=\psi(\varphi(0))$。我们于是有
\[
\begin{aligned}
\|\chi(a)-\chi(b)\|_{L^{p}}&\le\sup_{r\in[0,1]}\bigl\|(d\chi)_{ra+(1-r)b}\bigr\|^{\mathrm{op}}\,\|a-b\|_{L^{p}}\\
&\le\sup_{r\in[0,1]}\bigl\|(d\chi)_{ra+(1-r)b}\bigr\|^{\mathrm{op}}\,(C+K)\,\|Z\|_{W^{1,p}}\text{。}
\end{aligned}
\]
这里 $b=\psi(\varphi(0))=\psi(0)=(0,\widetilde{V}_{\rho}(0))$，而且我们知道\footnote{译注：下一式原书印作 $\widetilde{V}_{\rho}(s,t)=\dfrac{\partial\widetilde{J}}{\partial\lambda}(\lambda+\lambda_{\star},s+\rho)\cdot D_{1}\widetilde{\exp}_{(w_{\rho},Z)}\dfrac{\partial w_{\rho}}{\partial t}$；按其上文，$b$ 的第二个分量是 $\widetilde{V}_{\rho}(0)$，本式给出的应是它在 $(s,t)$ 处的值、且取 $Z=0$，故疑应作 $\widetilde{V}_{\rho}(0)(s,t)$，$\widetilde{\exp}$ 的下标亦应随之为 $(w_{\rho},0)$。此处径改。——译者注}
\[
\widetilde{V}_{\rho}(0)(s,t)=\frac{\partial\widetilde{J}}{\partial\lambda}(\lambda+\lambda_{\star},s+\rho)\cdot D_{1}\widetilde{\exp}_{(w_{\rho},0)}\,\frac{\partial w_{\rho}}{\partial t}\text{。}
\]
对一个取定的三元组 $(\lambda,s,t)$，这个向量的\trEuclidean{}范数被一个不依赖于 $\rho$ 的常数一致控制，后者显式地等于
\[
\sup_{(\lambda,s)\in[0,1]\times\RR}\Bigl\|\frac{\partial\widetilde{J}}{\partial\lambda}(\lambda,s)\Bigr\|\ \sup_{\substack{p\in W\\ \|Z\|\le r_{0}}}\bigl\|D_{1}\widetilde{\exp}_{(p,Z)}\bigr\|\ \sup_{(s,t)\in S^{1}\times\RR}\Bigl\|\frac{\partial w_{\rho}}{\partial t}\Bigr\|\text{。}
\]
我们在\cref{thm:dpsi-dphi-compose}的证明中继续，让 $\widetilde{V}_{\rho}$ 扮演 $Y_{0}^{\rho}$ 的角色，便得到\footnote{译注：此链尚不足以推出结论：均值不等式（本引理证明的第一行）所需者是 $\sup_{r\in[0,1]}\bigl\|(d\chi)_{ra+(1-r)b}\bigr\|^{\mathrm{op}}$ 本身，而这里（依\cref{thm:dpsi-dphi-compose}的条件 (3) 之类比）证得的只是它与 $\bigl\|(d\chi)_{b}\bigr\|^{\mathrm{op}}$ 之差，缺了基线一项；且第二行等号处原书还消失了一枚因子 $r$——$\bigl\|ra+(1-r)b-b\bigr\|_{L^{p}}=r\,\|a-b\|_{L^{p}}$，右端应为 $k_{3}r^{2}\,\|a-b\|_{L^{p}}$（或径略去首行多出的那个 $r$）。可照此补救：按上文 $b=\psi(0)=\bigl(0,\widetilde{V}_{\rho}(0)\bigr)$ 一致\trBounded，仿\cref{thm:dp-rho-estimate}条件 (4) 的验证可另取常数 $k_{7}>0$ 使 $\bigl\|(d\chi)_{b}\bigr\|^{\mathrm{op}}\le k_{7}$，从而 $\sup_{r\in[0,1]}\bigl\|(d\chi)_{ra+(1-r)b}\bigr\|^{\mathrm{op}}\le k_{7}+k_{3}(C+K)r_{0}$，最终常数即 $\bigl[k_{7}+k_{3}(C+K)r_{0}\bigr](C+K)$——引理的结论仍然成立，只是现下的推理链不成立。原书如此，照排。——译者注}
\[
\begin{aligned}
\bigl\|(d\chi)_{ra+(1-r)b}-(d\chi)_{b}\bigr\|^{\mathrm{op}}&\le k_{3}r\,\bigl\|ra+(1-r)b-b\bigr\|_{L^{p}}\\
&=k_{3}r\,\|a-b\|_{L^{p}}\\
&\le k_{3}(C+K)\,\|Z\|_{W^{1,p}}\\
&\le k_{3}(C+K)r_{0}\text{。}
\end{aligned}
\]
利用我们对 $\|\psi\varphi(Z)-\psi\varphi(0)\|$ 的估计，由此推出
\[
\Bigl\|\chi\psi\varphi(Z)-\chi\psi\varphi(0)\Bigr\|_{L^{p}}\le k_{3}(C+K)^{2}r_{0}\|Z\|_{W^{1,p}}\text{，}
\]
这就证明了\cref{thm:chi-psi-estimate}。
\end{proof}

我们现在继续建立\cref{thm:dfbar-lipschitz}的估计。由前一条引理（\cref{thm:chi-psi-estimate}），我们有
\[
\begin{aligned}
\bigl\|(d\overline{F}_{\rho})_{(\lambda,Z)}(a,Y)-(d\overline{F}_{\rho})_{(\lambda,0)}(a,Y)\bigr\|_{L^{p}}&\le k_{1}\,\|Z\|_{W^{1,p}}\|Y\|_{W^{1,p}}+|a|\,k_{2}\,\|Z\|_{W^{1,p}}\\
&=(k_{1}+k_{2})\left(|a|+\|Y\|_{W^{1,p}}\right)\|Z\|_{W^{1,p}}\text{，}
\end{aligned}
\]
从而
\[
\bigl\|(d\overline{F}_{\rho})_{(\lambda,Z)}-(d\overline{F}_{\rho})_{(\lambda,0)}\bigr\|^{\mathrm{op}}\le(k_{1}+k_{2})\,\|Z\|_{W^{1,p}}\text{。}
\]

我们就这样控制住了不等式右端的第一项：
\[
\bigl\|(d\overline{F}_{\rho})_{(\lambda,Z)}-(d\overline{F}_{\rho})_{(0,0)}\bigr\|^{\mathrm{op}}\le\bigl\|(d\overline{F}_{\rho})_{(\lambda,Z)}-(d\overline{F}_{\rho})_{(\lambda,0)}\bigr\|^{\mathrm{op}}+\bigl\|(d\overline{F}_{\rho})_{(\lambda,0)}-(d\overline{F}_{\rho})_{(0,0)}\bigr\|^{\mathrm{op}}\text{，}
\]
而我们尚须控制第二项。

对 $(a,Y)\in\RR\times W^{1,p}(\RR\times S^{1};\RR^{2n})$，我们写
\[
(d\overline{F}_{\rho})_{(\lambda,0)}(a,Y)=(dF_{\rho}^{\Gamma_{\lambda+\lambda_{\star}}})_{0}(Y)+a\bigl(V_{\rho}(\lambda+\lambda_{\star},0)\bigr)_{(Z_{i}^{\rho})}\text{，}
\]
其中
\[
V_{\rho}(\lambda+\lambda_{\star},0)=\frac{\partial J}{\partial\lambda}(\lambda+\lambda_{\star},s+\rho)\,\Bigl(\frac{\partial w_{\rho}}{\partial t}-X_{H_{s+\rho,t}^{\lambda+\lambda_{\star}}}\bigl(w_{\rho}\bigr)\Bigr)+\operatorname{grad}_{w_{\rho}}\frac{\partial H}{\partial\lambda}(\lambda+\lambda_{\star},s+\rho,t)\text{。}
\]
于是我们有
\[
\begin{aligned}
&\bigl\|(d\overline{F}_{\rho})_{(\lambda,0)}(a,Y)-(d\overline{F}_{\rho})_{(0,0)}(a,Y)\bigr\|_{L^{p}}\\
&\qquad\le\bigl\|(dF_{\rho}^{\Gamma_{\lambda+\lambda_{\star}}})_{0}(Y)-(dF_{\rho}^{\Gamma_{\lambda_{\star}}})_{0}(Y)\bigr\|_{L^{p}}\\
&\qquad\qquad+|a|\,\bigl\|\bigl(V_{\rho}(\lambda+\lambda_{\star},0)\bigr)_{Z_{i}^{\rho}}-\bigl(V_{\rho}(\lambda_{\star},0)\bigr)_{Z_{i}^{\rho}}\bigr\|_{L^{p}}\text{。}
\end{aligned}
\]
如同在\cref{subsec:construction-phi}中那样（并遵循\cref{thm:df-dbar-plus-s}的先例），
\[
(dF_{\rho}^{\Gamma_{\lambda+\lambda_{\star}}})_{0}(Y)=L_{\rho}^{\lambda}(Y)=\frac{\partial Y}{\partial s}+J_{0}\frac{\partial Y}{\partial t}+S(\lambda)Y\text{，}
\]
其中
\[
S:[0,1]\times\RR\times S^{1}\longrightarrow\mathrm{M}_{2n}(\RR)
\]
是一个 $\cC^{\infty}$ 类的映射。于是
\[
\begin{aligned}
\bigl\|(dF_{\rho}^{\Gamma_{\lambda+\lambda_{\star}}})_{0}(Y)-(dF_{\rho}^{\Gamma_{\lambda_{\star}}})_{0}(Y)\bigr\|_{L^{p}}&=\bigl\|S(\lambda)(Y)-S(0)(Y)\bigr\|_{L^{p}}\\
&\le k_{4}\,|\lambda|\,\|Y\|_{L^{p}}\\
&\le k_{4}\,|\lambda|\,\|Y\|_{W^{1,p}}\text{，}
\end{aligned}
\]
其中 $k_{4}=\sup_{\lambda\in[0,1]}\partial S/\partial\lambda$。其次，\footnote{译注：下一不等式左端原书印作 $\Bigl\|\bigl(V_{\rho}(\lambda+\lambda_{\star}),0-V_{\rho}(\lambda_{\star},0)\bigr)_{Z_{i}^{\rho}}\Bigr\|_{L^{p}}$；按前一段中已经出现的平行不等式 $\Bigl\|\bigl(V_{\rho}(\lambda+\lambda_{\star},0)\bigr)_{Z_{i}^{\rho}}-\bigl(V_{\rho}(\lambda_{\star},0)\bigr)_{Z_{i}^{\rho}}\Bigr\|_{L^{p}}$，第一个括号疑应作 $V_{\rho}(\lambda+\lambda_{\star},0)$。此处径改。——译者注}
\[
\begin{aligned}
&\Bigl\|\bigl(V_{\rho}(\lambda+\lambda_{\star},0)-V_{\rho}(\lambda_{\star},0)\bigr)_{Z_{i}^{\rho}}\Bigr\|_{L^{p}}\\
&\qquad\le\Bigl\|\Bigl(\frac{\partial J}{\partial\lambda}(\lambda+\lambda_{\star},s+\rho)-\frac{\partial J}{\partial\lambda}(\lambda_{\star},s+\rho)\Bigr)\Bigl(\frac{\partial w_{\rho}}{\partial t}\Bigr)_{Z_{i}^{\rho}}\Bigr\|_{L^{p}}\\
&\qquad\qquad+\Bigl\|\Bigl(\operatorname{grad}_{w_{\rho}}\Bigl(\frac{\partial H}{\partial\lambda}(\lambda+\lambda_{\star},s+\rho,t)-\frac{\partial H}{\partial\lambda}(\lambda_{\star},s+\rho,t)\Bigr)\Bigr)_{Z_{i}^{\rho}}\Bigr\|_{L^{p}}\\
&\qquad\le\Bigl(\sup_{\substack{\lambda\in[0,1]\\ s\in\RR}}\Bigl\|\frac{\partial^{2}J}{\partial\lambda^{2}}(\lambda,s)\Bigr\|\Bigr)|\lambda|\,\Bigl\|\Bigl(\frac{\partial w_{\rho}}{\partial t}\Bigr)_{Z_{i}^{\rho}}\Bigr\|_{L^{p}}\\
&\qquad\qquad+|\lambda|\!\!\sup_{\substack{\lambda\in[0,1]\\ (s,t)\in\RR\times S^{1}}}\Bigl\|\Bigl(\operatorname{grad}_{w_{\rho}}\frac{\partial^{2}H}{\partial\lambda^{2}}(\lambda,s,t)\Bigr)_{Z_{i}^{\rho}}\Bigr\|_{L^{p}}\le k_{5}\,|\lambda|\text{。}
\end{aligned}
\]
这就导致
\[
\bigl\|(d\overline{F}_{\rho})_{(\lambda,0)}(a,Y)-(d\overline{F}_{\rho})_{(0,0)}(a,Y)\bigr\|_{L^{p}}\le k_{4}\,|\lambda|\,\|Y\|_{W^{1,p}}+k_{5}\,|a|\,|\lambda|\text{，}
\]
这表明
\[
\bigl\|(d\overline{F}_{\rho})_{(\lambda,0)}-(d\overline{F}_{\rho})_{(0,0)}\bigr\|^{\mathrm{op}}\le k_{6}\,|\lambda|\text{。}
\]
把这个不等式同先前得到的
\[
\bigl\|(d\overline{F}_{\rho})_{(\lambda,Z)}-(d\overline{F}_{\rho})_{(\lambda,0)}\bigr\|^{\mathrm{op}}\le(k_{1}+k_{2})\,\|Z\|_{W^{1,p}}\text{，}
\]
结合起来，我们最终发现
\[
\begin{aligned}
\bigl\|(d\overline{F}_{\rho})_{(\lambda,Z)}-(d\overline{F}_{\rho})_{(0,0)}\bigr\|^{\mathrm{op}}&\le(k_{1}+k_{2})\,\|Z\|_{W^{1,p}}+k_{6}\,|\lambda|\\
&\le k\bigl(\|Z\|_{W^{1,p}}+|\lambda|\bigr)\text{，}
\end{aligned}
\]
这正是我们结束\cref{thm:dfbar-lipschitz}的证明所需要的。
\end{proof}

在\cref{subsec:properties-psi-rho}中，为了建立 $\psi$ 的那些所需要的性质，我们使用过下面这条陈述，如今正是证明它的时候了。

\emph{引理（\cref{thm:dfbar-drho-estimate}）。}存在一个不依赖于 $\rho$ 的常数 $C>0$，使得对 $\rho\ge\rho_{0}$，我们有
\[
\Bigl\|\frac{\partial\overline{F}_{\rho}}{\partial\rho}\bigl(\overline{\lambda}_{\rho},\gamma_{\rho}\bigr)-\frac{\partial\overline{F}_{\rho}}{\partial\rho}(0,0)\Bigr\|_{L^{p}}\le C\,\bigl\|(\overline{\lambda}_{\rho},\gamma_{\rho})\bigr\|_{\RR\times W^{1,p}}\text{。}
\]

\begin{proof}
我们将\cref{thm:df-drho-estimate}应用于诸算子 $F^{\Gamma_{\lambda}}$（$\lambda\in[0,1]$）（参见\cref{thm:gamma-newton-picard}的证明）。于是，对 $Z\in\overline{B}(0,r_{0})$，我们有
\[
\Bigl\|\frac{\partial F^{\Gamma_{\lambda}}}{\partial\rho}(\rho,Z)-\frac{\partial F^{\Gamma_{\lambda}}}{\partial\rho}(\rho,0)\Bigr\|_{L^{p}}\le M\|Z\|_{W^{1,p}}\text{，}
\]
其中 $M>0$ 是一个不依赖于 $\rho$ 的常数（也不依赖于 $\lambda$，因为 $\lambda\in[0,1]$）。特别地，由于 $\overline{F}_{\rho}(\lambda,Z)=F_{\rho}^{\Gamma_{\lambda+\lambda_{\star}}}(Z)$，我们有
\[
\Bigl\|\frac{\partial\overline{F}_{\rho}}{\partial\rho}\bigl(\overline{\lambda}_{\rho},\gamma_{\rho}\bigr)-\frac{\partial\overline{F}_{\rho}}{\partial\rho}\bigl(\overline{\lambda}_{\rho},0\bigr)\Bigr\|_{L^{p}}\le M\|\gamma_{\rho}\|_{W^{1,p}}\text{。}
\]
为了结束这条引理的证明，只需证明存在一个不依赖于 $\rho$ 的常数 $C\ge 0$，使得
\[
\Bigl\|\frac{\partial\overline{F}_{\rho}}{\partial\rho}\bigl(\overline{\lambda}_{\rho},0\bigr)-\frac{\partial\overline{F}_{\rho}}{\partial\rho}(0,0)\Bigr\|_{L^{p}}\le C\,\bigl|\overline{\lambda}_{\rho}\bigr|\text{。}
\]
对任意的 $\lambda\in\RR$，我们有
\[
\begin{aligned}
\overline{F}_{\rho}(\lambda,0)&=\Bigl(\scr{F}_{\rho}^{\Gamma_{\lambda+\lambda_{\star}}}(w_{\rho})\Bigr)_{Z_{i}^{\rho}}\\
&=\Bigl(\frac{\partial w_{\rho}}{\partial s}+J_{s+\rho}^{\lambda+\lambda_{\star}}(w_{\rho})\frac{\partial w_{\rho}}{\partial t}+\operatorname{grad}H_{s+\rho,t}^{\lambda+\lambda_{\star}}(w_{\rho})\Bigr)_{Z_{i}^{\rho}}\text{。}
\end{aligned}
\]
设 $\chi(\rho,\lambda)=\overline{F}_{\rho}(\lambda,0)\in L^{p}(\RR\times S^{1};\RR^{2n})$。作为 $(s,t)$ 的函数，$\dfrac{\partial\chi}{\partial\lambda}(\rho,\lambda)$ 具有紧\trSupport。事实上，$J_{s+\rho}^{\lambda+\lambda_{\star}}$ 与 $\operatorname{grad}H_{s+\rho,t}^{\lambda+\lambda_{\star}}$ 对 $\lambda$ 的导数在 $|s+\rho|\ge R$ 时为零，因为诸\trIsotopy{} $\Gamma_{\lambda}$ 在 $|s|\ge R$ 时都是常值的。于是，作为 $(s,t)$ 的函数，$\dfrac{\partial^{2}\chi}{\partial\rho\partial\lambda}$ 也具有紧\trSupport。从而
\[
\begin{aligned}
\Bigl\|\frac{\partial\chi}{\partial\rho}(\rho,\lambda)-\frac{\partial\chi}{\partial\rho}(\rho,0)\Bigr\|_{L^{p}}&\le|\lambda|\cdot\Bigl\|\sup_{\lambda}\frac{\partial^{2}\chi}{\partial\rho\partial\lambda}\Bigr\|_{L^{p}}\\
&=|\lambda|\cdot\Bigl\|\sup_{\lambda}\frac{\partial^{2}\chi}{\partial\rho\partial\lambda}\Bigr\|_{L^{p}([-R-\rho,R-\rho]\times S^{1};\RR^{2n})}\\
&\le K\,|\lambda|\cdot\sup_{\substack{\lambda\in\RR\\ \rho\ge\rho_{0}\\ (s,t)\in[-R-\rho,R-\rho]\times S^{1}}}\Bigl\|\frac{\partial^{2}\chi}{\partial\rho\partial\lambda}(s,t)\Bigr\|
\end{aligned}
\]
（最后一项中的范数是 $\RR^{2n}$ 的\trEuclidean{}范数）。与 $\partial w_{\rho}/\partial t$ 的范数一样，诸 $Z_{i}^{\rho}$ 的范数以及它们对 $\rho$ 的导数都被一些不依赖于 $\rho$ 的常数所控制。对 $J_{s+\rho}^{\lambda+\lambda_{\star}}$ 与 $\operatorname{grad}H_{s+\rho,t}^{\lambda+\lambda_{\star}}$ 关于 $\rho$ 与 $\lambda$ 的导数也同样如此。于是，$\partial^{2}\chi/\partial\rho\partial\lambda$ 对一切 $(\lambda,\rho,s,t)\in\RR\times[\rho_{0},+\infty[\,\times\,\RR\times S^{1}$（按\trEuclidean{}范数）都\trBounded，这就结束了\cref{thm:dfbar-drho-estimate}的证明。
\end{proof}
